\PassOptionsToPackage{draft}{hyperref}
\documentclass[12pt,a4paper]{article}

% Compilar con LuaLaTeX + Biber:  .\build.ps1
\usepackage{viu-mrob-thesis}
\usepackage{csquotes}
\usepackage[
  backend=biber,
  style=apa,
  natbib=true,
  hyperref=false,
  sorting=nyt,
  sortcites=true
]{biblatex}
\addbibresource{references.bib}
\renewcommand*{\bibfont}{\raggedright}

\usepackage{array}
\usepackage{tabularx}
\usepackage{longtable}
\usepackage{booktabs}
\usepackage{needspace}

\setviuheadertitle{Revisión de literatura}
\hypersetup{
  pdftitle={Revisión de literatura: coordinación multi-robot, coaliciones y transporte cooperativo},
  pdfauthor={Jorge Luis Mayorga Taborda},
  pdfsubject={Estado del arte para el TFM VIU-MROB 2026}
}

% Marca de trabajo pendiente. Se resalta en el PDF para que ningún hueco
% llegue por descuido a una entrega. Antes de congelar el subdocumento no
% puede quedar ninguna: `grep -rn "\pendiente" sections/`.
\newcommand{\pendiente}[1]{%
  \par\noindent\colorbox{VIUOrange!12}{%
    \parbox{\dimexpr\linewidth-2\fboxsep\relax}{%
      \footnotesize\color{VIUDark}\textbf{[PENDIENTE]}\ #1}}\par\smallskip}

\begin{document}

\begin{center}
  {\LARGE\bfseries Revisión de literatura\par}
  \vspace{2mm}
  {\large Coordinación multi-robot, formación de coaliciones\
   y transporte cooperativo\par}
  \vspace{4mm}
  {\normalsize Jorge Luis Mayorga Taborda\par}
  {\footnotesize Máster Universitario en Robótica --- VIU\par}
  {\footnotesize Subdocumento de trabajo. No es un capítulo cerrado.\par}
\end{center}

\vspace{4mm}
\tableofcontents
\vspace{6mm}

\input{sections/01-alcance-y-metodo}
\input{sections/02-asignacion-de-tareas}
\input{sections/03-formacion-de-coaliciones}
\input{sections/04-consenso-y-control-distribuido}
\input{sections/05-juegos-y-dinamicas-poblacionales}
\input{sections/06-transporte-cooperativo}
\input{sections/07-sintesis-y-hueco}

\printbibliography[title={Referencias}]

\end{document}
